\subsection{VII)}

{\large\begin{align*}\displaystyle F(A, B, C, D) &= \prod(1, 3, 5, 7, 13, 15)\end{align*}}

\subsubsection{Obtener la mínima expresión booleana}
Esto es una productoria, por lo que se trata de maxitérminos y cada posición representa $0's$ de la función, en un producto de sumas. Sí hay un 1 en una entrada, hay que negarla.

\begin{center}
\begin{tabular}{c|cccc|c}
\rowcolor[HTML]{FFCCC9} 
\cellcolor[HTML]{9AFF99}M & $A$ & $B$ & $C$ & $D$ & \cellcolor[HTML]{FFFFC7}$F(A,B,C,D) $\\ \hline
0                         & 0   & 0   & 0   & 0   &                                  \\
1                         & 0   & 0   & 0   & 1   & 0                                \\
2                         & 0   & 0   & 1   & 0   &                                  \\
3                         & 0   & 0   & 1   & 1   & 0                                \\
4                         & 0   & 1   & 0   & 0   &                                  \\
5                         & 0   & 1   & 0   & 1   & 0                                \\
6                         & 0   & 1   & 1   & 0   &                                  \\
7                         & 0   & 1   & 1   & 1   & 0                                \\
8                         & 1   & 0   & 0   & 0   &                                  \\
9                         & 1   & 0   & 0   & 1   &                                  \\
10                        & 1   & 0   & 1   & 0   &                                  \\
11                        & 1   & 0   & 1   & 1   &                                  \\
12                        & 1   & 1   & 0   & 0   &                                  \\
13                        & 1   & 1   & 0   & 1   & 0                                \\
14                        & 1   & 1   & 1   & 0   &                                  \\
15                        & 1   & 1   & 1   & 1   & 0                               
\end{tabular}
\end{center}
{\large
\begin{align*}
    F(A,B,C,D) &= (A+B+C+\overline{D})\cdot{}(A+B+\overline{C}+\overline{D})\cdot{}(A+\overline{B}+C+\overline{D})\cdot{}(A+\overline{B}+\overline{C}+\overline{D})
\end{align*}
    \begin{align*}
(\overline{A}+\overline{B}+C+\overline{D})\cdot{}(\overline{A}+\overline{B}+\overline{C}+\overline{D})
    \end{align*}
}
\begin{center}
    \begin{karnaugh-map}[4][4][1][$C\qquad D$][$A\qquad B$]
        \maxterms{1,3,5,7,13,15}
        \implicant{1}{7}
        \implicant{5}{15}
    \end{karnaugh-map}
\end{center}

\begin{itemize}[nosep]
    \item Primer bloque ($M_1$, $M_3$, $M_5$ y $M_7$) cambia $B$ y $C$, va $A$ y $D$ negado.
    \item Segundo bloque ($M_5$,$M_7$,$M_13$ y $M_15$) cambia $A$ Y $C$, va $B$ negado y $D$ negado.
    \item La función simplificada queda:
\end{itemize}

{\large
    \begin{align*}
        F(A,B,C,D) &= (A+\overline{D})\cdot{}(\overline{B}+\overline{D})
    \end{align*}
}

\noindent Se puede sacar factor común $\overline{D}$:
{\large
    \begin{align*}
        F(A,B,C,D) &= \overline{D} + A\cdot{}\overline{B}
    \end{align*}
}

\subsubsection{Simulación}

\imagen[]{10cm}{./imagenes/ejercicio7Falstad.png}\subsubsection{Descripción en Verilog}\begin{lstlisting}
module descripcionVerilogEjercicio7(
    input inputA,
    input inputB,
    input inputC,
    input inputD,
    output outputFunction
    );
	 assign outputFunction = ~inputD | inputA & (~inputB);


endmodule
     \end{lstlisting}\subsubsection{RTL output}\imagen[Bloque RTL]{15cm}{./imagenes/rtlEjercicio7Verilog.png}\imagen[Implementación con compuertas]{15cm}{./imagenes/rtlCompuertasEjercicio7Verilog.png}\subsubsection{Simulación temporal}\noindent Código del módulo Verilog Test Fixture:\begin{lstlisting}

module verilogTestFixtureEjercicio7;

	// Inputs
	reg inputA;
	reg inputB;
	reg inputC;
	reg inputD;

	// Outputs
	wire outputFunction;

	// Instantiate the Unit Under Test (UUT)
	descripcionVerilogEjercicio7 uut (
		.inputA(inputA), 
		.inputB(inputB), 
		.inputC(inputC), 
		.inputD(inputD), 
		.outputFunction(outputFunction)
	);

	initial begin
		// Initialize Inputs
		inputA = 0;
		inputB = 0;
		inputC = 0;
		inputD = 0;

		// Wait 100 ns for global reset to finish
		#100;
        
		// Add stimulus here
		inputA = 0; inputB = 0; inputC = 0; inputD = 1;
		#100;
		inputA = 0; inputB = 0; inputC = 1; inputD = 0;
		#100;
		inputA = 0; inputB = 0; inputC = 1; inputD = 1;
		#100;
		inputA = 0; inputB = 1; inputC = 0; inputD = 0;
		#100;
		inputA = 0; inputB = 1; inputC = 0; inputD = 1;
		#100;
		inputA = 0; inputB = 1; inputC = 1; inputD = 0;
		#100;
		inputA = 0; inputB = 1; inputC = 1; inputD = 1;
		#100;
		inputA = 1; inputB = 0; inputC = 0; inputD = 0;
		#100;
		inputA = 1; inputB = 0; inputC = 0; inputD = 1;
		#100;
		inputA = 1; inputB = 0; inputC = 1; inputD = 0;
		#100;
		inputA = 1; inputB = 0; inputC = 1; inputD = 1;
		#100;
		inputA = 1; inputB = 1; inputC = 0; inputD = 0;
		#100;
		inputA = 1; inputB = 1; inputC = 1; inputD = 0;
		#100;
		inputA = 1; inputB = 1; inputC = 1; inputD = 1;

	end
      
endmodule
 \end{lstlisting}\imagen[Simulación temporal]{15cm}{./imagenes/simulacionTemporalEjercicio7.png}
